\documentclass[10pt,a4paper]{article}
\usepackage[margin=1.5cm]{geometry}
\usepackage{amsmath,graphicx,booktabs,enumitem}
\setlist{nosep,leftmargin=*}
\pagestyle{empty}
\setlength{\parindent}{0pt}
\setlength{\parskip}{3pt}
\newcommand{\sect}[1]{\par\vspace{2pt}{\large\textbf{#1}}\par}

\begin{document}

{\Large\textbf{Feasibility Review: AD vs CN Classification on ADNI with ConvNeXt}}\\
COMP3710 Pattern Recognition \quad\textbar\quad William O'Driscoll, 48849470

\vspace{0.3cm}
\sect{1. User need, scope and acceptance criteria}
\textbf{User:} A clinician reviewing screening T1-weighted brain MRIs who needs a classifier that decides confident cases automatically and routes uncertain ones to an expert. \textbf{Scope.} A prototype classifying 2D ADNI slices as Alzheimer's disease (AD) or cognitively normal (CN), and testing whether its confidence is trustworthy compared with simpler baselines (the open dilemma). Targets are provisional.
\begin{enumerate}
\vspace{0.1cm}
  \item \textbf{Accuracy:} per-slice test accuracy $\geq 0.80$ on a subject-level split, with a 95\% bootstrap interval resampled over subjects.
  \vspace{0.1cm}
  \item \textbf{Beyond the shortcut:} exceed the slice-window-only baseline (66.2\% on test) and, within each mixed-class window, exceed that window's majority-class rate.
  \vspace{0.1cm}
  \item \textbf{Calibration:} ECE and reliability diagrams for the MLP, CNN and ConvNeXt, comparing confidence on correct and wrong predictions; ConvNeXt compared with the baselines on macro-F1, AUROC and ECE.
  \vspace{0.1cm}
  \item \textbf{Reject option:} a confidence threshold giving $\geq 0.90$ accuracy on accepted slices at $\geq 50\%$ coverage, with coverage and residual errors reported.
  \vspace{0.1cm}
  \item \textbf{Resources:} peak GPU memory $\leq 16$~GB and $\leq 6$~h per training run on one A100.
\end{enumerate}

\vspace{0.3cm}
\sect{2. Model choice and course concepts}
\textbf{Data:} ADNI AD/CN 2D slices (256$\times$240 grayscale, 20 per scan), zero-padded to 256$\times$256 so stage resolutions halve cleanly. \textbf{Model:} a reduced ConvNeXt (Liu et al., 2022; Hard difficulty) trained from scratch, as pretrained weights are not permitted. ConvNeXt-T uses four stages with 3, 3, 9 and 3 blocks and widths of 96 to 768. The reduced model keeps the same block design (4$\times$4 patchify stem, 7$\times$7 depthwise convolution, LayerNorm, 4$\times$ inverted bottleneck with a single GELU, layer scale and stochastic depth) but uses fewer blocks and channels to suit the dataset size and compute budget. In the paper's ablations these design steps raise ImageNet-1K top-1 accuracy from 78.8\% to 82.0\% (ResNet-50 regime). That evidence comes from far larger data, so transfer to small brain-MRI data is tested, not assumed. Training uses AdamW with weight decay and a cosine schedule. Augmentation is mild (small rotation, shift, scale and intensity jitter) with no flips, because left--right flips of these sagittal-looking slices would swap the front and back of the head.
\textbf{Baselines:} window-only, MLP and CNN on identical splits, giving a progression from no anatomy, to no spatial structure, to local filters, to a modern design. \textbf{Evaluation:} per slice, as recommended for 2D models, with confidence intervals resampled over subjects. \textbf{Course concepts:} convolutional inductive bias, normalisation, regularisation (weight decay, stochastic depth, augmentation), probabilistic calibration and selective prediction (risk--coverage).

\vspace{0.3cm}
\sect{3. Preliminary feasibility evidence}
\textbf{Data audit} (1{,}526 scans, 680 subjects, 30{,}520 slices): the provided train/test folders share 216 of 331 test subjects, so all scans are pooled and re-split by subject (70/15/15, seed 42). The test set holds 102 subjects (33 AD), 219 scans and 4{,}380 slices. No subject has both classes. Slice window is strongly tied to class: a window-only rule fitted on the training split scores 66.2\% on test (65.5\% validation) against 55.7\% for always predicting the training majority, so about 10 points of any result can come from slice depth alone. Accuracy is therefore also reported per window. \textbf{Pipeline test:} a short run (3 epochs on a 2{,}000-slice subset) exercised data loading, the CNN, the training loop and checkpointing end to end; accuracy at that scale is not meaningful.

\vspace{0.3cm}
\sect{4. Risks, budget, next experiment and fallback}
\textbf{Risks and mitigations:} (i) the subject-level split lowers accuracy relative to the leaking provided split, and 0.80 may be hard from scratch; stochastic depth, weight decay, augmentation and early stopping on validation address overfitting. (ii) The slice-window shortcut; the window-only baseline, per-window accuracy and a window check on failure cases expose it. (iii) A small effective test set (102 subjects) gives wide intervals; differences between models are claimed only where bootstrap intervals separate. (iv) GPU queue waits of about 12~h limit the number of runs; jobs are submitted early as batch scripts, tested on small subsets first, with several configurations per job where possible. \textbf{Budget:} one A100 through Slurm batch jobs, each capped at 6~h, with a planned total of about eight runs (CNN, MLP, three to four ConvNeXt configurations and final seeds). \textbf{Next experiment:} a full CNN baseline run, then a first ConvNeXt run after a one-batch overfitting check. \textbf{Fallback:} if ConvNeXt misses 0.80, first reduce width and depth or strengthen regularisation; if it is still short, report the result, since the calibration, reject-option and failure analysis carry the dilemma investigation. A second architecture (GFNet) is only attempted after the core analysis is complete.

\end{document}